%% Shared page design for manuscript.tex and supplementary.tex (plain article class; Nature has no required LaTeX class
%% and takes a PDF at first submission). Two modes:
%%   reading mode (default)       : single-spaced, no line numbers -- a clean copy to read and circulate;
%%   submission mode              : double-spaced with line numbers, as the Nature formatting guide asks for a PDF submission.
%% Submission mode is switched on by defining \submissionmode before this file is read (see manuscript_submission.tex).

\usepackage[a4paper,left=22mm,right=22mm,top=24mm,bottom=26mm]{geometry}
\usepackage{iftex}
\ifPDFTeX                                   % pdfLaTeX (Overleaf default)
  \usepackage[T1]{fontenc}
  \usepackage{newtxtext,newtxmath}          % Times-like text and matching mathematics
  \usepackage[scaled=0.9]{helvet}           % Helvetica-like sans serif for title, headings, captions and tables
\else                                       % XeLaTeX / LuaLaTeX / tectonic
  \usepackage{newtxtext,newtxmath}
  \setsansfont{texgyreheros}[Extension=.otf,UprightFont=*-regular,BoldFont=*-bold,ItalicFont=*-italic,BoldItalicFont=*-bolditalic,Scale=0.9]
\fi
\usepackage{graphicx}
\usepackage{xcolor}
\usepackage{amsmath}
\usepackage{booktabs,tabularx,array}
\usepackage{float}
\usepackage{setspace}
\usepackage{url}
\usepackage[super,comma,sort&compress]{natbib}   % Nature: superscript reference numbers
\usepackage{titlesec}
\usepackage{caption}
\usepackage[hidelinks]{hyperref}

\bibliographystyle{sn-nature}
\providecommand{\bibcommenthead}{}          % written by sn-nature.bst; defined by the Springer class only
\renewcommand{\bibnumfmt}[1]{#1.}           % Nature reference list: "1." not "[1]"
\setlength{\bibsep}{2pt}
\renewcommand{\bibfont}{\small}

%% Headings: bold sans serif, unnumbered, close to the text they introduce.
\titleformat{\section}{\sffamily\bfseries\large\raggedright}{}{0pt}{}
\titlespacing*{\section}{0pt}{16pt plus 2pt minus 2pt}{4pt}
\titleformat{\subsection}{\sffamily\bfseries\normalsize\raggedright}{}{0pt}{}
\titlespacing*{\subsection}{0pt}{10pt plus 2pt minus 1pt}{2pt}
\newcommand{\bmhead}[1]{\subsection*{#1}}

%% Captions in Nature style: "Fig. 1 | Title. Text" and "Table 1 | Title", small sans serif, full width.
\DeclareCaptionLabelSeparator{bar}{\ \textbar\ }
\captionsetup{font={footnotesize,sf},labelfont=bf,labelsep=bar,justification=justified,singlelinecheck=false,skip=6pt}
\captionsetup[table]{position=top,textfont=bf,skip=4pt}
\renewcommand{\figurename}{Fig.}

%% One font for every table (the table files call \tablefont).
\newcommand{\tablefont}{\sffamily\fontsize{7.5}{9.5}\selectfont}
\renewcommand{\arraystretch}{1.1}

%% Floats: allow large figures and tables at the top of a page without pushing them to the end.
\renewcommand{\topfraction}{0.9}
\renewcommand{\bottomfraction}{0.6}
\renewcommand{\textfraction}{0.08}
\renewcommand{\floatpagefraction}{0.75}
\setlength{\textfloatsep}{14pt plus 2pt minus 2pt}
\floatplacement{table}{t}
\floatplacement{figure}{t}

\setlength{\parindent}{1.2em}
\setlength{\parskip}{0pt}
\emergencystretch=2em                       % no line may stick out into the margin
\frenchspacing
\raggedbottom

\ifdefined\submissionmode
  \usepackage{lineno}
  \renewcommand{\linenumberfont}{\normalfont\scriptsize\sffamily\color{black!55}}
  \setlength{\linenumbersep}{8pt}
  \AtBeginDocument{\doublespacing\linenumbers}
\fi

\newcommand{\todo}[1]{\textcolor{red}{[TODO: #1]}}
\newcommand{\muB}{\ensuremath{\mu_\mathrm{B}}}
\newcommand{\Tc}{\ensuremath{T_\mathrm{C}}}
\newcommand{\mcell}{\ensuremath{m_\mathrm{cell}}}
\newcommand{\mfu}{\ensuremath{m_\mathrm{fu}}}
